\section{算法设计与实现}
教学模型直接返回解析解 $x^*=2$，再计算基线点 $1$ 与解析解的目标值。
程序不靠有限误差采样证明最坏界，公式已由前节证明。
真实问题应说明输入输出、步骤、停止条件、失败处理和复杂度。

\begin{table}[htbp]
\centering
\caption{教学示例的理论与实现映射}
\begin{tabularx}{\linewidth}{lX}
\toprule 内容 & 对应入口 \\
\midrule
名义及稳健目标 & \texttt{examples/demo.py} 的 \texttt{objectives} \\
演示表格生成 & 同文件 \texttt{main}，输出到 \texttt{generated/} \\
论文构建 & \texttt{build.py}，双遍编译正文和AI说明 \\
\bottomrule
\end{tabularx}
\end{table}
